% Reader-first front matter for the corrected Version 2 manuscript.
% This file is not included automatically.
%
% IMPORTANT STATUS BOUNDARY
% -------------------------
% The post-closure abstract and theorem below must not be inserted into main.tex
% until the global Section VIII attained-demand/all-deficit physical-fibre
% theorem and exact aggregate weight identity have been checked.

% ---------------------------------------------------------------------------
% A. AUDIT-SAFE ABSTRACT FOR THE CURRENT RESEARCH BRANCH
% ---------------------------------------------------------------------------

\begin{abstract}
For a graph $G$, let $\chi(G)$ be its chromatic number and let $\zeta(G)$ be the
least number of parts in a vertex partition in which every part induces either
an empty graph or a complete graph.  We develop a full-sequence candidate proof
that, for $G_n\sim G(n,1/2)$, the difference $\chi(G_n)-\zeta(G_n)$ is bounded
below with high probability by a positive constant times
$n/(\log n)^3$.  The argument has three layers.  An ordinary first-moment
calculation locates the chromatic threshold; a signed four-size first moment
places the cocolouring threshold lower by order $n/(\log n)^3$; and an exact
signed-overlap identity reduces the required second moment to partial
diagonals, a matching of high cells, and a residual even-subgraph attachment.
The corrected Section VIII route sums high cells through literal physical
matching fibres, one global falling-factorial denominator, and a geometric
all-deficit estimate.  The residual attachment is then controlled by injective
restriction outside the exposed matching.  A rare signed seed is amplified to
high probability by leftover colouring and bounded differences.  The remaining
submission-blocking step is the global reindexing of every attained canonical
high skeleton by block support, admissible deficits, and local partial physical
matchings, with exact aggregate weight preservation.  Until that finite theorem
is closed, the normalized second moment and the final random-graph theorem
remain conditional.
\end{abstract}

% ---------------------------------------------------------------------------
% B. POST-CLOSURE PUBLICATION ABSTRACT
% Delete the audit-safe abstract above and use this version only after closure.
% ---------------------------------------------------------------------------

\iffalse
\begin{abstract}
For a graph $G$, let $\chi(G)$ be its chromatic number and let $\zeta(G)$ be the
least number of parts in a vertex partition in which every part induces either
an empty graph or a complete graph.  For $G_n\sim G(n,1/2)$ we prove, uniformly
over the full independence-threshold phase, that
\[
  \chi(G_n)-\zeta(G_n)
  \ge
  \left[
    \frac{(\log2)^2}{8}
      \bigl(\log2-D_4(\delta_n)\bigr)-o(1)
  \right]
  \frac{n}{(\log n)^3}
\]
with high probability.  In particular,
\[
  \chi(G_n)-\zeta(G_n)
  \ge
  \left[
    \frac{(\log2)^2}{8}
      \log\!\left(\frac{1000}{639}\right)-o(1)
  \right]
  \frac{n}{(\log n)^3}
\]
with high probability.  The gain comes from signed colour classes: at density
$1/2$, declaring a class independent or complete has the same edge cost and
contributes an exact factor $2$ per class.  An exact overlap identity separates
local cell weights from a binary cycle-space factor.  We sum the canonical high
cells through literal physical matching fibres and a single global
falling-factorial denominator, and control the residual even-subgraph family by
injective restriction outside the exposed matching.  The resulting second
moment gives a positive-probability signed seed, which is amplified to high
probability by leftover colouring and bounded differences.
\end{abstract}
\fi

% ---------------------------------------------------------------------------
% C. READER-FIRST INTRODUCTION
% The theorem box is written in post-closure form; retain the explicit status
% paragraph while this remains a research draft.
% ---------------------------------------------------------------------------

\section*{Introduction}\label{introduction-reader-first-v2}

Let $G_n\sim G(n,1/2)$ be the random graph on $[n]$.  The chromatic number
$\chi(G)$ is the least number of independent sets in a partition of $V(G)$.
The cochromatic number $\zeta(G)$ is the least number of parts in a partition in
which every part is either independent or complete.  Erd\H{o}s and Gimbel asked
whether
\[
  \chi(G_n)-\zeta(G_n)\longrightarrow\infty
\]
with high probability.  The question is now catalogued as Erd\H{o}s
Problem~625.

The mechanism behind the gap is simple.  Fix a partition into $k$ classes.  An
ordinary colouring requires every class to be independent.  A cocolouring may
declare each class independent or complete.  Because $G(n,1/2)$ is invariant
under complementation, the two declarations have the same probability cost.
Summing over the sign choices therefore contributes an exact factor $2^k$.
That additional entropy lowers the signed first-moment root by order
$n/(\log n)^3$.

The proof has three logically distinct layers.

\begin{enumerate}
\item \emph{Location.}  We compare the ordinary-colouring root $r_+(n)$ with
      the signed four-size root $r_4^{\mathrm{co}}(n)$.
\item \emph{Existence.}  We prove that the lower signed root is populated by
      controlling one normalized second moment.
\item \emph{Amplification.}  We turn the resulting positive-probability signed
      seed into a high-probability cocolouring without changing the leading
      root separation.
\end{enumerate}

Writing
\[
  A_4(\delta):=\log2-D_4(\delta),
\]
the phase-uniform first-moment comparison gives
\[
  r_+(n)-r_4^{\mathrm{co}}(n)
  =
  \left[
    \frac{(\log2)^2}{4}A_4(\delta_n)+o(1)
  \right]
  \frac{n}{(\log n)^3}.
\]
We place the signed seed at the midpoint between these roots, retaining one
half of the separation.  The exact entropy certificate gives the uniform bound
\[
  A_4(\delta)>\log\!\left(\frac{1000}{639}\right).
\]

% Use the following theorem box only after the Section VIII closure theorem.
\begin{resultbox}{Theorem 1: conditional Version 2 target}
For $G_n\sim G(n,1/2)$,
\[
  \chi(G_n)-\zeta(G_n)
  \ge
  \left[
    \frac{(\log2)^2}{8}A_4(\delta_n)-o(1)
  \right]
  \frac{n}{(\log n)^3}
\]
with high probability.  Consequently,
\[
  \chi(G_n)-\zeta(G_n)
  \ge
  \left[
    \frac{(\log2)^2}{8}
      \log\!\left(\frac{1000}{639}\right)-o(1)
  \right]
  \frac{n}{(\log n)^3}
\]
with high probability.
\end{resultbox}

\displayheading{What the second moment must prove}
Let $Z_n$ count the signed four-size profiles at the selected midpoint.  The
entire overlap analysis is directed toward the single estimate
\[
  \frac{\mathbb E Z_n^2}{(\mathbb E Z_n)^2}
  \le
  \exp\!\left\{
    o\!\left(\frac{n}{(\log n)^4}\right)
  \right\}.
\]
The second-moment sections do not move either first-moment root.  Their purpose
is only to show that the signed first-moment advantage survives overlap.

For two ordered signed partitions, let $r=(r_{ab})$ be the overlap table and
let $H_r$ be the bipartite support graph of cells with multiplicity at least
two.  Summing the sign declarations gives the exact identity
\[
  \operatorname{SignedOverlapWeight}(r)
  =
  \left(\prod_{a,b}g(r_{ab})\right)2^{\beta(H_r)},
\]
where $\beta(H_r)$ is the binary cycle rank.  The local product records the cell
multiplicities, while the cycle-space factor records the remaining sign
compatibility.

We then reduce the overlap in three stages.  Common whole classes form the
partial diagonals.  The remaining canonical high cells have support on a block
matching.  Everything else is a residual even-subgraph attachment after the
high matching has been exposed.

\displayheading{The high-cell dictionary}
A \emph{profile block} is an actual class; a \emph{block atom} is its type/slot
label.  The \emph{block support} $P$ is the matching of block pairs carrying
positive high demand.  For $e\in P$, let $s_e,t_e$ be the endpoint sizes,
\[
  m_e:=\min\{s_e,t_e\},
  \qquad
  j_e:=m_e-h_e,
\]
where $j_e$ is the actual high multiplicity and $h_e$ is its deficit from full
containment.  A \emph{partial physical matching} is the literal size-$j_e$
matching of stubs in one selected block pair.  An \emph{endpoint table} records
only how many support edges have each pair of endpoint types.  A partial
physical matching need not have a unique full completion; the proof compares
whole finite fibres.

For fixed support $P$ and multiplicities $j=(j_e)$, the aggregate physical
weight is
\[
  w(P,j)
  =
  \frac{\displaystyle
    \prod_{e\in P}(s_e)_{j_e}(t_e)_{j_e}}
  {\displaystyle
    (n)_J\prod_{e\in P}j_e!}
  \prod_{e\in P}g(j_e),
  \qquad
  J:=\sum_{e\in P}j_e.
\]
There is one factorial $j_e!$ per selected cell and one global denominator
$(n)_J$.

If $H:=\sum_e h_e$, then the full-reference denominator changes by
\[
  \frac{(n)_{J+H}}{(n)_J}=(n-J)_H\le n^H.
\]
For endpoint sizes $m,m+d$, define
\[
  R_{m,d}(h)
  :=
  \frac{\binom mh}{(d+1)(d+2)\cdots(d+h)}
  2^{-hm+h(h+1)/2}.
\]
The global comparison is
\[
  \frac{w(P,m-h)}{w_{\mathrm{full}}(P)}
  \le
  \prod_{e\in P}n^{h_e}R_{m_e,d_e}(h_e).
\]
The high-cell condition implies $2h_e<m_e$, turning the sum over all deficits
into a geometric product.  Combined with endpoint transport, this yields
\[
  \operatorname{BareSkeletonSum}_n
  \le
  \exp\!\left\{
    o\!\left(\frac{n}{(\log n)^4}\right)
  \right\}.
\]

\displayheading{Residual attachment}
Let $M$ be the exposed high-support matching.  Restriction outside $M$ is
injective on even residual edge sets: the symmetric difference of two
completions would be an even subset of a matching and hence empty.  Therefore,
for nonnegative activities $q_e$,
\[
  \sum_{F\ \mathrm{even}}
    \prod_{e\in F\setminus M}q_e
  \le
  \prod_{e\notin M}(1+q_e).
\]
The local increment activity is bounded by the same $q$ activity.  The two
residual regimes then give
\[
  \operatorname{AttachmentSum}_n
  \le
  \operatorname{BareSkeletonSum}_n
  \exp\!\left\{
    \eta_n\frac{n}{(\log n)^4}
  \right\},
  \qquad \eta_n\to0.
\]
This completes the normalized second moment once the physical-fibre reindexing
has been established.

\displayheading{Amplification and final intersection}
Paley--Zygmund gives a signed seed with positive probability.  A leftover
colouring lemma and a one-Lipschitz cocolourable-capacity variable, combined
with bounded differences, promote this seed to a high-probability upper bound
for $\zeta(G_n)$ with lower-order loss.  The ordinary first-moment analysis gives
the corresponding high-probability lower bound for $\chi(G_n)$.  The two events
need not be independent; their intersection follows from a union bound.

\displayheading{Current audit boundary}
At the time of this draft, the remaining submission-blocking theorem is the
global identification of every attained canonical high skeleton with its block
support, admissible deficits, and local partial physical matching fibres,
together with exact aggregate weight preservation.  The theorem box and the
post-closure abstract above must remain labelled conditional until that finite
assembly has passed the focused proof audit.
